\section{Recommendations}

\subsection{Scaling the portfolio to 100 papers}
The 60-paper cohort (Wave~1: 30 papers in $\sim$3 weeks; Wave~2: 15 papers in 5 days; Wave~3: 10 papers in 1 day; Wave~4: 5 papers in 1 day) demonstrates approximately 3$\times$ throughput improvement with subagent parallelism and domain-specific API integration. Wave~1 took approximately three weeks of agent wall-clock,
including deferred waiting on HPC allocations. With 60 papers now complete,
we see a super-linear payoff toward 100 if two
investments are made:
\begin{itemize}
\item \textbf{Replication-plan cache.} Every paper now has a
\texttt{replication\_plan\_<osti>.tex}; treating these as a searchable corpus
lets the agent reuse environment recipes and avoid re-solving ``how do I build
AMBER+OpenMM on Aurora'' for each new paper.
\item \textbf{Domain pre-staged environments.} Containers for the top 6 domain
buckets (materials DFT, MSM/OpenMM, graph-algorithms, power/Simulink,
bioinformatics/BV-BRC, combustion LES) would cut agent setup time from hours
to minutes.
\end{itemize}
At 100 papers, we expect total elapsed time on the order of six weeks of agent
wall-clock, most of which is queue latency rather than agent reasoning.

\subsection{Compute that would unblock the bottom quartile}
\begin{itemize}
\item A Polaris allocation with $\sim 5\times 10^4$ node-hours would unblock
1275503 (CROC), 1559043 (full Fig.\,3 ignition ensemble), and 2396968 v2
(6-band LSST latent SDE).
\item Dedicated 8$\times$A100 for two weeks would unblock 3003857
(full Kolmogorov) and 1868518 (8500-node grid DRL).
\item BV-BRC SEEDtk access (request sent 2026-04-24) unblocks two computational
biology replications simultaneously.
\end{itemize}

\subsection{Where AI-agent replication excels, and where it falters}
\textbf{Excels.} Closed-loop numerical experiments with open code + public
data; paper sections where the pipeline is well-specified; tasks where
the agent can iterate against a unit-test-style scalar metric. The
combinatorial, graph-algorithmic, and topological-response papers in our
cohort were completed to 90--100\% fidelity.

\textbf{Falters.} Experimental-only papers where the agent cannot access the
instrument; proprietary / gated codebases; papers whose claims depend on
long-tail HPC campaigns the agent cannot afford; papers where the key
contribution is tacit (``we chose force field X because'') and not written down.

\subsection{Process improvements (cohort v2)}
\begin{enumerate}
\item \textbf{Pre-registration.} Before running, the agent should commit a
scored coverage \emph{target} per paper. This converts the agreement score
from a retrospective assessment to a calibration signal.
\item \textbf{Seed budgets.} Require $\geq 3$ stochastic seeds for any DL/RL
claim, logged with wall-clock.
\item \textbf{External validators.} Wire in xTB / xtb-python / DFT surrogates
as standard filters so chemistry papers don't fall back to RDKit
approximations silently.
\item \textbf{Author contact protocol.} Build a light-touch
``clarification email'' tool, gated on human approval, for the $\sim 10\%$ of
papers where a single question to the author would close a 2-point coverage
gap.
\item \textbf{Score calibration.} Audit a random 10\% of scorings against a
second reviewer; we expect $\pm 1$ on each axis for well-calibrated scoring.
\end{enumerate}


